\documentclass[11pt]{article}

\usepackage[margin=1in]{geometry}
\usepackage{amsmath,amssymb,amsthm,mathtools}
\usepackage{booktabs}
\usepackage[hidelinks]{hyperref}
\usepackage[nameinlink,capitalize]{cleveref}

\newtheorem{theorem}{Theorem}[section]
\newtheorem{lemma}[theorem]{Lemma}
\newtheorem{proposition}[theorem]{Proposition}
\newtheorem{corollary}[theorem]{Corollary}
\theoremstyle{definition}
\newtheorem{definition}[theorem]{Definition}
\theoremstyle{remark}
\newtheorem{remark}[theorem]{Remark}

\newcommand{\cO}{\mathcal O}
\newcommand{\cM}{\mathfrak M}
\newcommand{\F}{\mathbb F}
\newcommand{\Q}{\mathbb Q}
\newcommand{\R}{\mathbb R}
\newcommand{\C}{\mathbb C}
\newcommand{\rd}{\operatorname{rd}}
\newcommand{\Nm}{\operatorname N}
\newcommand{\Cl}{\operatorname{Cl}}

\title{Distance-Sidon subsets from bounded-inertia number-field towers%
\thanks{This investigation was carried out with substantial
assistance from automated systems, and we omit a conventional byline rather
than overstate the named researchers' personal contributions.
\emph{Nikol Savova}, University of Oxford,
\href{mailto:nikol.p.savova@gmail.com}{\texttt{nikol.p.savova@gmail.com}};
\emph{Sihao Huang}, independent researcher,
\href{mailto:sihao.c.huang@gmail.com}{\texttt{sihao.c.huang@gmail.com}}.}}

\author{}
\date{\today}

\begin{document}
\maketitle

\begin{abstract}
For a finite planar point set $P$, let $s(P)$ be the largest size of a
subset whose pairwise distances are all distinct, and put
\[
 F_2(n)=\min_{|P|=n}s(P).
\]
Recent work gives the window
$n^{1/3}\ll F_2(n)\ll n^{1/2-\varepsilon}$ for some absolute
$\varepsilon>0$.  We make the upper-bound saving explicit:
\[
 \boxed{F_2(n)\ll n^{0.49368323}}.
\]
The construction combines three ingredients.  We impose bounded tame
inertia and bounded Frobenius order in a Golod--Shafarevich pro-$2$ tower;
we place points in a product of disks in the Eisenstein CM extension of a
totally real layer; and we amplify repeated algebraic norms using
prime-power valuation patterns.  A one-sided ideal-packing argument gives
the geometric constant $2\sqrt3/\pi$.  The tower is built over
$\Q(\sqrt{11235917})$.  Its class, ray, Frobenius, and numerical endpoint
data are finite and checked by three executable implementations, including
an independent reconstruction.
\end{abstract}

\section{Introduction}

Call a finite set $A\subset\R^2$ \emph{distance-Sidon} if the map
\[
 \{x,y\}\longmapsto |x-y|,
 \qquad x,y\in A,\quad x\ne y,
\]
is injective.  For a finite planar set $P$, write
\[
 s(P)=\max\{|A|:A\subseteq P\text{ is distance-Sidon}\},
 \qquad F_2(n)=\min_{|P|=n}s(P).
\]
The problem of estimating $F_2(n)$ is recorded as Erd\H{o}s problem
1208~\cite{ErdosProblems1208}.  Charalambides proved
$F_2(n)\gg n^{1/3}/\log n$~\cite{Charalambides2013}; Clemen, F\"uhrer,
and Roche-Newton recently removed the logarithm~\cite{ClemenEtAl2026}.
In the other direction, Lee, Pohoata, and Zhu constructed point sets for
which every sufficiently large subset contains a repeated
distance~\cite{LeePohoataZhu2026}.  In particular, for some absolute
$\varepsilon>0$,
\[
 n^{1/3}\ll F_2(n)\ll n^{1/2-\varepsilon}.       \tag{1.1}
\]
Before this polynomial saving, Croot, Mao, Pohoata, Sheffer, and Yip obtained
the first super-polylogarithmic improvement for distance-Sidon subsets of the
square grid~\cite{CrootEtAl2026}.

Our purpose is to give a compact quantitative version of the upper-bound
construction.

\begin{theorem}\label{thm:main}
For every positive integer $n$, there is an $n$-point set
$P\subset\R^2$ such that every distance-Sidon subset of $P$ has size
$O(n^{0.49368323})$.  Equivalently,
\[
 F_2(n)\ll n^{0.49368323}.
\]
\end{theorem}

The proof is an arithmetic amplification of the Minkowski-grid idea.  Its
architecture is easier to describe than the numerical exponent suggests.
An infinite totally real pro-$2$ tower supplies fields of dyadic degree and
bounded root discriminant.  At many unramified primes, Frobenius has order
at most two.  After adjoining a primitive cube root of unity, those primes
split, so a difference $z$ admits several prescribed allocations of its
valuations between the conjugate factors of $z\bar z$.  If a subset had all
distances distinct, these allocations could not occur too often.  Comparing
a lower bound from congruence classes with an upper bound for totally
positive ideal elements yields the master inequality in
\cref{prop:master}.

Two refinements account for the explicit saving.  First, forcing the square
of each tame inertia generator to be trivial limits every ramification
index to two while preserving the Frattini quotient.  Second, a product of
planar disks in an Eisenstein CM field has a better diameter-to-volume ratio
than the coordinate boxes used in the initial construction.  The remaining
optimization is a one-dimensional concave knapsack problem.

The bounded-inertia input goes back to Hajir and
Maire~\cite{HajirMaire2001,HajirMaire2002}, building on Koch's class-tower
criterion~\cite{Hajir1996,Koch2002}; we use the later tower-cutting framework
of Hajir, Maire, and Ramakrishna~\cite{HajirMaireRamakrishna2021}.  The recent
number-field construction for the unit-distance problem and Sawin's explicit
version are direct methodological precedents~\cite{AlonEtAl2026,Sawin2026}.
N\"aslund's subsequent ideal-lattice covolume refinement is the closest
geometric analogue of our product-disk step~\cite{Naslund2026}.  We state the
tame, totally real Shafarevich presentation theorem in the precise form
needed below.  The field-specific calculations use PARI/GP~\cite{PARI2}; the
accompanying verifier reconstructs the relevant Kummer kernel, checks every
selected prime, and evaluates the final inequalities at two independent
precisions.

All logarithms below are natural.

\section{The Eisenstein norm sieve}

Let $L$ be a totally real field of degree $m$, and put
\[
 K=L(\zeta_3),\qquad \zeta_3^2+\zeta_3+1=0.
\]
Choose one complex embedding of $K$ above each real embedding of $L$.
This embeds $\cO_K$ as a lattice in $\C^m$.

\subsection{A product of disks}

For $R>0$, let $\Delta_R\subset\C$ be a disk of area $R^2$, and set
$\mathcal D_R=\Delta_R^m$.  The complex Minkowski lattice has covolume
\[
 \operatorname{covol}(\cO_K)=2^{-m}\sqrt{|\operatorname{disc}K|}.
\]
The relative discriminant of $L(\zeta_3)/L$ divides $3\cO_L$, whence
\[
 \operatorname{covol}(\cO_K)
 \le \left(\frac{\sqrt3}{2}\rd(L)\right)^m.       \tag{2.1}
\]
Averaging translates over a fundamental domain therefore proves the next
elementary fact.

\begin{lemma}[disk averaging]\label{lem:disk}
There is a translate of $\mathcal D_R$ containing at least
\[
 \left(\frac{2R^2}{\sqrt3\,\rd(L)}\right)^m
\]
elements of $\cO_K$.  If $z$ is the difference of two elements in that
translate and $\eta=z\bar z\in\cO_L$, then, for every real embedding
$\sigma:L\hookrightarrow\R$,
\[
 0<\sigma(\eta)\le \frac4\pi R^2.                \tag{2.2}\label{eq:norm-box}
\]
\end{lemma}

\begin{proof}
The first assertion is volume divided by covolume.  A disk of area $R^2$
has diameter $2R/\sqrt\pi$.  Under the complex embedding above $\sigma$,
$\sigma(z\bar z)=|z|^2$, which proves \eqref{eq:norm-box}.
\end{proof}

We shall use
\[
 R^2=\frac{\sqrt3}{2}\rd(L)n^{1/m}.              \tag{2.3}\label{eq:R-choice}
\]
Then the selected translate contains at least $n$ lattice points.  After
projection through one distinguished complex embedding, it gives a planar
point set.  That projection is injective.  Moreover, equality of two
projected squared distances implies equality of the corresponding elements
$z\bar z\in L$, because a field embedding is injective.

The following one-sided estimate is the reason that total positivity is
useful.

\begin{lemma}[one-sided ideal packing]\label{lem:ideal-packing}
Let $\mathfrak a\subset\cO_L$ be a nonzero integral ideal and $Y>0$.
Then
\[
 \#\{\eta\in\mathfrak a:0<\sigma(\eta)\le Y
       \text{ for every }\sigma:L\hookrightarrow\R\}
 \le \left(1+\frac{Y}{\Nm(\mathfrak a)^{1/m}}\right)^m. \tag{2.4}
\]
\end{lemma}

\begin{proof}
Partition $(0,Y]^m$ into half-open cubes of side
$\delta=\Nm(\mathfrak a)^{1/m}$.  Two embedded ideal elements in one cube
would have a nonzero difference $\gamma\in\mathfrak a$ with
$|\Nm_{L/\Q}\gamma|<\delta^m=\Nm\mathfrak a$.  This is impossible because
$(\gamma)\subseteq\mathfrak a$.  There are at most
$(1+Y/\delta)^m$ cubes.
\end{proof}

\subsection{Valuation patterns}

Suppose a prime ideal $\mathfrak q$ of $L$ splits in $K/L$ as
\[
 \mathfrak q\cO_K=\mathfrak P\bar{\mathfrak P}.
\]
Fix a depth $k\ge0$.  For each $0\le a\le k$, consider the additive
lattice
\[
 I_a=\mathfrak P^a\bar{\mathfrak P}^{k-a}\subset\cO_K. \tag{2.5}
\]
It has index $\Nm(\mathfrak q)^k$.  Products over several prime ideals and
possibly different depths give pattern lattices $I_{\mathbf a}$.  Write
\[
 \cM=\prod_{\mathfrak q}\mathfrak q^{k_{\mathfrak q}},
 \quad
 H=\prod_{\mathfrak q}(k_{\mathfrak q}+1),
 \quad
 \Lambda=\prod_{\mathfrak q}
     \left(1+\Nm\mathfrak q^{-1}+\cdots+
                 \Nm\mathfrak q^{-k_{\mathfrak q}}\right),             \tag{2.6}
\]
and $\mu=\Nm(\cM)^{1/m}$.

For a finite set $A\subset\cO_K$, Cauchy--Schwarz over the cosets of every
pattern lattice gives
\[
 \sum_{\mathbf a}\#\{(x,y)\in A^2:x\ne y, x-y\in I_{\mathbf a}\}
 \ge H\frac{|A|(|A|-\Nm\cM)}{\Nm\cM}.           \tag{2.7}\label{eq:pair-lower}
\]
On the other hand, if $z=x-y$ belongs to several patterns at a fixed prime,
the excess multiplicity forces an extra power of $\mathfrak q$ to divide
$z\bar z$.  More precisely, the number of patterns containing $z$ is at
most the number of ideals $\mathfrak b\mid\cM$ for which
\[
 \cM\mathfrak b\mid(z\bar z).                    \tag{2.8}\label{eq:divisor-switch}
\]
This is the divisor switch: at $\mathfrak P$ and
$\bar{\mathfrak P}$, the admissible values of $a$ form the interval
$[k-v_{\bar{\mathfrak P}}(z),v_{\mathfrak P}(z)]\cap[0,k]$.

If the planar projection of $A$ is distance-Sidon, each nonzero value
$z\bar z$ comes from at most two ordered pairs.  Applying
\cref{lem:ideal-packing} to \eqref{eq:divisor-switch}, and then summing
$\Nm(\mathfrak b)^{-1}$ over $\mathfrak b\mid\cM$, yields the master bound.

\begin{proposition}[norm-sieve master inequality]\label{prop:master}
Let $A$ be a distance-Sidon subset of a translate of $\mathcal D_R$.
With the notation above,
\[
 |A|\le \mu^m+\sqrt2R^m
 \left[(\Lambda/H)^{1/m}
       \left(\frac4\pi+\frac{\mu^2}{R^2}\right)
 \right]^{m/2}.                                  \tag{2.9}\label{eq:master}
\]
\end{proposition}

\begin{proof}
By \eqref{eq:norm-box}, every norm $z\bar z$ lies in the positive box with
$Y=4R^2/\pi$.  For $\mathfrak b\mid\cM$,
\cref{lem:ideal-packing} gives
\[
 \#\{\eta\in\cM\mathfrak b:0<\sigma(\eta)\le Y\}
 \le\left[
 \frac{R^2}{\mu\Nm(\mathfrak b)^{1/m}}
 \left(\frac4\pi+\frac{\mu^2}{R^2}\right)
 \right]^m,
\]
where we used $\Nm(\mathfrak b)^{1/m}\le\mu$.  Sum over
$\mathfrak b\mid\cM$, use
$\sum_{\mathfrak b\mid\cM}\Nm(\mathfrak b)^{-1}=\Lambda$, multiply by
the ordered-pair factor two, and compare with \eqref{eq:pair-lower}.  If
$|A|\ge\mu^m$, then
$(|A|-\mu^m)^2\le |A|(|A|-\mu^m)$; taking square roots gives
\eqref{eq:master}.
The case $|A|<\mu^m$ is immediate.
\end{proof}

Combining \eqref{eq:R-choice} with \eqref{eq:master} replaces the real
two-coordinate
constant $4/\pi$ by the effective constant
\[
 C_{\rm Eis}=\frac{2\sqrt3}{\pi}.                 \tag{2.10}
\]

\section{A bounded-inertia pro-$2$ tower}

We record the group-theoretic input in a form that separates the standard
presentation theorem from the two elementary quotient operations used here.

For a finite set $T$ of odd prime ideals of a number field $E$, put
\[
 V_T=\{a\in E^\times:(a)=\mathfrak a^2\text{ for some fractional ideal }
 \mathfrak a,\ a\in E_{\mathfrak p}^{\times2}\text{ for every }
 \mathfrak p\in T\}/E^{\times2}.                  \tag{3.1}
\]
This is the usual tame Kummer obstruction group.

\begin{theorem}[tame totally real Shafarevich--Koch presentation]
\label{thm:shafarevich}
Let $E$ be totally real with $r_1$ real places, let $T$ be a finite set of
odd prime ideals, and let $G_T^+$ be the Galois group of the maximal pro-$2$
extension of $E$ unramified away from $T$ in which every real place splits.
Then $G_T^+$ has a finite minimal presentation with generator and relation
ranks satisfying
\[
 d(G_T^+)=|T|-r_1+\dim_{\F_2}V_T,
 \qquad
 r(G_T^+)\le |T|-1+\dim_{\F_2}V_T.               \tag{3.2}
\]
In particular, if $E$ is real quadratic and $V_T=0$, then
\[
 d(G_T^+)=|T|-2,\qquad r(G_T^+)\le |T|-1=d(G_T^+)+1. \tag{3.3}\label{eq:shaf-bound}
\]
\end{theorem}

\begin{proof}
This is the tame $p=2$ presentation theorem of
Koch~\cite[Theorems~11.5 and~11.8]{Koch2002}; see also
\cite[Section~10.7]{NSW2008}.  In the standard formulas,
\[
 d=|T|-(r_1+r_2)+1-\delta_2(E)+\dim_{\F_2}V_T,
 \qquad
 r\le |T|-\delta_2(E)+\dim_{\F_2}V_T.
\]
Here $r_2=0$ and $\delta_2(E)=1$ because $-1\in E$.  The dyadic places lie
outside $T$, while requiring all real places to split is the totally real
local condition; this is the archimedean refinement that removes the extra
$p=2$ relation present for the unrestricted extension.  Substitution gives
\eqref{eq:shaf-bound}.
\end{proof}

\begin{proposition}[tower criterion]\label{prop:tower}
Let $E$ be totally real, and let $T$ be a finite set of odd prime ideals of
$E$.  Suppose the maximal totally real pro-$2$ extension unramified away
from $T$ has a presentation with generator rank $d$ and relation rank at
most $r_0$.  Let $U$ be a set of unramified prime ideals with chosen
Frobenius elements.  Form the quotient obtained by imposing
\[
 \tau_{\mathfrak p}^2=1\quad(\mathfrak p\in T),
 \qquad
 \operatorname{Frob}_{\mathfrak q}^2=1\quad(\mathfrak q\in U),         \tag{3.4}\label{eq:square-caps}
\]
where $\tau_{\mathfrak p}$ generates tame pro-$2$ inertia.  If
\[
 4(r_0+|T|+|U|)<d^2,                              \tag{3.5}\label{eq:gs-condition}
\]
then this quotient is infinite.  Its Frattini quotient is unchanged, and
in every finite layer its ramification index at $\mathfrak p\in T$ is at
most two.
\end{proposition}

\begin{proof}
Tame pro-$2$ inertia at an odd prime is procyclic.  Normal closure of the
single relation $\tau_{\mathfrak p}^2$ therefore caps all conjugate inertia
groups above $\mathfrak p$; normal closure likewise makes the chosen
Frobenius representative irrelevant.  Every relation in
\eqref{eq:square-caps} is a square and hence lies in the Frattini subgroup,
so the generator rank remains $d$.
Adding at most $|T|+|U|$ relations gives the relation bound in
\eqref{eq:gs-condition}.  The strict Golod--Shafarevich inequality then proves
infinitude; see, for example, \cite{Ershov2012}.  The ramification assertion
is immediate from the inertia cap.
\end{proof}

In the explicit field below, both $V_T=0$ and the generator space are
verified exactly.  Thus \cref{thm:shafarevich}, rather than a generic rank
estimate, supplies the base relation bound.

At a tame prime $\mathfrak p$ of $E$, the relative different exponent in a
finite layer is at most one half of the relative degree.  Thus the tower in
\cref{prop:tower} has root discriminant bounded by
\[
 D_0=\rd(E)\prod_{\mathfrak p\in T}
              \Nm_E(\mathfrak p)^{1/(2[E:\Q])}.         \tag{3.6}
\]
For a real quadratic base, the exponent is $1/4$.

We also need the Frobenius caps to make primes split in the Eisenstein CM
extension of each layer.  Call an unramified prime ideal $\mathfrak q$
\emph{useful} if it does not lie above $3$ and, writing
$Q=\Nm_E\mathfrak q$, either $Q\equiv1\pmod3$, or $Q\equiv2\pmod3$ and the
Frobenius class of $\mathfrak q$ is nonzero in the Frattini quotient.  In the
second case it has exact order two after the square cap, and the residue
field in the selected layers has size $Q^2\equiv1\pmod3$; in the first case
splitting is automatic.

For completeness, an infinite finitely generated pro-$2$ group has, for
every sufficiently large $j$, an open normal subgroup $N_j$ contained in
its Frattini subgroup and of index $2^j$.  Indeed, take arbitrarily large
finite $2$-group quotients below the Frattini subgroup and refine a normal
series by central factors of order two.  The corresponding Galois layers
have every sufficiently large dyadic degree and preserve every nonzero
Frattini class.

\section{The local frontier}

Let a prime ideal of the quadratic base have norm $Q$.  Across an orbit in
a tower layer, one extra unit of valuation depth has absolute-degree
normalized cost
\[
 c_Q=\frac12\log Q.                                \tag{4.1}\label{eq:item-cost}
\]
For $0<t<1$ and $k\ge1$, put
\[
 A_k(t)=\frac{k+1}{k}\frac{1-t^k}{1-t^{k+1}}.     \tag{4.2}
\]
The guaranteed gain of the $k$th depth increment is
\[
 g_{Q,k}=\frac14\log A_k(Q^{-2}).                 \tag{4.3}\label{eq:item-gain}
\]
This is exact when the residue degree is two and is a lower bound when it
is one.

\begin{lemma}[all-depth comparison]\label{lem:all-depth}
For $0<t<1$ and $k\ge1$,
\[
 A_k(t)^2\ge A_k(t^2).
\]
Moreover $A_k(t)$ decreases with $k$.
\end{lemma}

\begin{proof}
After cancelling positive factors, the first assertion is equivalent to
\[
 (1-t)\left(\sum_{j=0}^{2k}t^j-(2k+1)t^k\right)\ge0,
\]
which is AM--GM.  For the second, write
\[
 b_k(t)=\frac{1+t+\cdots+t^k}{k+1}
       =\frac1{1-t}\int_t^1x^k\,dx.
\]
Cauchy--Schwarz makes the moment sequence $b_k(t)$ log-convex, while
$A_k(t)=b_{k-1}(t)/b_k(t)$.
\end{proof}

Sort all available items $(c_Q,g_{Q,k})$ by decreasing slope
$g_{Q,k}/c_Q$, subject to taking the depths at each prime in order.  Let
$F(x)$ be the fractional-knapsack gain at total cost $x$.  By
\cref{lem:all-depth}, $F$ is increasing, concave, and realizable in high
degree up to an error $O(1/m)$: assign the final fractional item to the
largest collection of prime ideals whose total residue degree remains below
the desired fraction.

Put
\[
 w=\frac{\log n}{2m},\qquad x=\log\mu.
\]
Set $x=2\alpha w$, so the first term in \eqref{eq:master} is $n^\alpha$.
Using \eqref{eq:R-choice}, the second term is at most a constant times
$n^\alpha$
provided
\[
 F(2\alpha w)\ge
 \log(C_{\rm Eis}D_0)+(2-4\alpha)w+
 \log\left(1+
 \frac{e^{2(2\alpha-1)w}}{C_{\rm Eis}D_0}\right).       \tag{4.4}\label{eq:frontier}
\]
The difference between the two sides is concave in $w$: the first term is
concave, the middle term is affine, and
$-\log(1+ae^{-cw})$ is concave for $a,c>0$.  Consequently it is enough to
check \eqref{eq:frontier} at the two endpoints of an interval
$[w_0,2w_0]$.

\begin{lemma}[phase criterion]\label{lem:phase}
Suppose the tower has layers of every sufficiently large dyadic degree and
\eqref{eq:frontier} holds with a fixed positive margin at $w=w_0$ and
$w=2w_0$.
Then $F_2(n)\ll n^\alpha$.
\end{lemma}

\begin{proof}
For every sufficiently large $n$, choose a tower layer of degree $m$ such
that $w=(\log n)/(2m)$ lies in $[w_0,2w_0)$.  Concavity gives
\eqref{eq:frontier}; the fixed endpoint margin absorbs the $O(1/m)$ placewise
rounding error.  Apply \cref{prop:master} to the $n$ ambient points supplied
by \cref{lem:disk}.  The finitely many smaller values of $n$ are absorbed
by the implied constant.
\end{proof}

\section{The explicit tower and certificate}

We now specialize to
\[
 E=\Q(\sqrt D),\qquad D=11235917=7\cdot11\cdot337\cdot433. \tag{5.1}
\]
The ring of integers has basis $1,\omega$, where
\[
 \omega^2-\omega-2808979=0.
\]
Let $T$ be the first $217$ odd prime ideals of $E$, ordered first by norm
and then by the underlying rational prime; ties between the two split ideals
above one rational prime are broken by the increasing residue of $\omega$.

\begin{proposition}[finite arithmetic certificate]\label{prop:certificate}
For the field and prime set above, the following statements hold.
\begin{enumerate}
\item
$\Cl(E)\cong C_{14}\times C_2$ and
$\Cl^+(E)\cong C_{14}\times C_2\times C_2$.
The $T$-class group is trivial.
\item
The space of squareclasses $aE^{\times2}$ with $(a)$ an ideal square has
dimension four.  Its $4\times217$ matrix of local quadratic characters at
$T$ has rank four.  Consequently $V_T=0$.
\item
The $219$ fundamental $T$-unit squareclasses have rank-four image under the
two sign and modulo-$4$ conditions.  The safe Kummer kernel therefore has
dimension $d=215$.
\item
There are $11123$ useful prime ideals after $T$, with no rejection before
the cutoff.  The last selected and useful ideals have data
\[
 (1063,1063,963),\qquad (121951,121951,70091),      \tag{5.2}
\]
where the triples record norm, underlying rational prime, and the residue
of $\omega$.
\item
The tower quotient has relation rank at most
\[
 r=216+217+11123=11556,
 \qquad 4r=46224=215^2-1<215^2.                   \tag{5.3}\label{eq:gs-certificate}
\]
Its root discriminant satisfies
\[
 \log D_0=\frac12\log D+\frac14\sum_{\mathfrak p\in T}
                      \log\Nm\mathfrak p
 =318.9527585916460414427765165\ldots.             \tag{5.4}
\]
\end{enumerate}
\end{proposition}

The class and ray assertions in \cref{prop:certificate} are certified by
PARI/GP with \texttt{bnfcertify}=1, so they are unconditional and use no
GRH assumption.  They are not used as opaque numerical output: the verifier expands
the returned $T$-units in the basis $1,\omega$, recomputes both real signs
and the square classes modulo $4$, and independently obtains the same
rank-four row space.  It also constructs the four global ideal-square
classes and verifies rank four of their local-character matrix, proving
$V_T=0$.  For every candidate norm congruent to two modulo
three, it evaluates the quadratic-character functional on this exact
$215$-dimensional kernel.  All $11123$ required ideals pass.

The numerical part uses the rigorous rational estimate
\[
 C_{\rm Eis}=\frac{2\sqrt3}{\pi}<\frac{71603}{64935}. \tag{5.5}\label{eq:safe-constant}
\]
It constructs the sorted frontier from
\eqref{eq:item-cost}--\eqref{eq:item-gain}.  At
\[
 \alpha=0.49368323,\qquad
 w_0=39928.004292781350011665354078\ldots,         \tag{5.6}
\]
the margins in \eqref{eq:frontier} at $w_0$ and $2w_0$ are respectively
\[
 0.001632817036815572250219059\ldots,
 \qquad
 0.003247075212037706995029833\ldots.              \tag{5.7}\label{eq:margins}
\]
The active terminal slopes are approximately $0.0306672$ and $0.0191268$.
The largest omitted fourth-depth slope is
\[
 0.009999973433845534938505168\ldots,               \tag{5.8}\label{eq:omitted-slope}
\]
so \cref{lem:all-depth} excludes every omitted later depth.  The two
fixed-anchor threshold crossings lie below $0.49368323$; equalizing them
gives
\[
 \alpha_*=0.4936832199308881880151091959\ldots.     \tag{5.9}
\]
Thus the advertised exponent clears the equalized threshold by only
$1.0069\ldots\times10^{-8}$.  The relation budget in
\eqref{eq:gs-certificate} is likewise deliberately saturated at the largest
integer allowed by the strict quadratic Golod--Shafarevich inequality:
$4r=d^2-1$.

As insurance against weakening the presentation interface, the verifier
also recomputes desaturated certificates.  If one assumes only
$r_0\le d+2$, omitting the final useful prime gives threshold
$0.4936833032744\ldots$ and hence the safe bound
$F_2(n)\ll n^{0.4936834}$.  Under $r_0\le d+6$, omitting the final five
useful primes gives threshold $0.4936836367746\ldots$ and the safe bound
$F_2(n)\ll n^{0.4936837}$.

\begin{proof}[Proof of \cref{thm:main}]
By \cref{prop:certificate} and the tame Shafarevich bound
\eqref{eq:shaf-bound}, the strict inequality \eqref{eq:gs-certificate}
allows \cref{prop:tower} to produce an
infinite totally real tower with root discriminant at most $D_0$ and all
selected useful primes of residue degree at most two.  The Frobenius test
ensures that those primes split after passing to the Eisenstein CM
extension.  Equations
\eqref{eq:safe-constant}--\eqref{eq:omitted-slope} verify the two hypotheses of
\cref{lem:phase}.  The lemma gives
$F_2(n)\ll n^{0.49368323}$.
\end{proof}

\section{Verification and scope}

The proof has a conceptual part and a finite certificate.  The conceptual
part is contained in the lemmas and propositions of Sections~2--4; none of
those statements depends on the particular quadratic field.  The certificate
checks:
\begin{itemize}
\item the BNF, ordinary and narrow class groups, and localized class group;
\item the four-dimensional ideal-square space and the equality $V_T=0$;
\item the exact $T$-unit basis and its sign/modulo-$4$ kernel;
\item every selected prime ideal and every Frobenius usefulness test;
\item the integer Golod--Shafarevich inequality, both desaturated fallbacks,
and the root discriminant;
\item the complete active depth frontier, omitted-depth cutoff, and both
endpoint inequalities at $100$ and $150$ decimal digits.
\end{itemize}
Each endpoint computation is stable on increasing the working precision from
$100$ to $150$ decimal digits, while the advertised margins exceed $10^{-3}$.
The repository includes a second hostile verifier which independently
reconstructs the Kummer row space and repeats the arithmetic and endpoint
checks; its output agrees far beyond the displayed precision.

\paragraph{Data and code availability.}
The source package accompanying this manuscript contains a dependency-closed,
single-field verifier, its expected output, and the independent hostile
reconstruction.  The public repository is
\url{https://github.com/NikolSavova/proof_hunter}.  The immutable certificate
snapshot is the
\href{https://github.com/NikolSavova/proof_hunter/tree/c6f68c4b8dc458242f317f2c5f3d76ce4df36e5f/phase2/loop/erdos1208/paper/certificate}%
{repository at commit \texttt{c6f68c4}}.

This paper does not determine the order of $F_2(n)$.  The gap between
$1/3$ and $0.49368323$ remains substantial.  Nor do we claim that
$\Q(\sqrt{11235917})$ is globally optimal for the construction.  The point
of the explicit field is to make the quantitative theorem finite and
reproducible.  A significant approach toward the conjectural cube-root
scale will require a new global mechanism rather than further decimal
optimization of the same tower.

\bibliographystyle{alpha}
\bibliography{references}

\end{document}
